\batchmode


\documentclass[a4paper,11pt,twoside,openright]{article}
\RequirePackage{ifthen}



\usepackage[french]{babel}
\usepackage[latin1]{inputenc}

\usepackage{subfigure}      

\usepackage{indentfirst}
\usepackage{minitoc}
\usepackage{url}
\usepackage{algorithmic}
\usepackage{algorithm}
\usepackage{graphicx}
\usepackage{fancyhdr}
\usepackage{amsmath}
\usepackage{amssymb}
\usepackage{amsfonts}
\usepackage{amsthm}
\usepackage{multirow}
\usepackage{vmargin}
%
\providecommand{\comment}[1]{}
\comment{


 
\floatstyle{boxed}

 
\restylefloat{figure}
\begin{figure}[htbp]
{\small \begin{displaymath}
\begin{array}{ll}
    \rotatebox[origin=c]{90}{Programme} \left\{
          \begin{array}{ll}
                \rotatebox[origin=c]{0}{Description} & \left\{
        \begin{array}{ll}
        \rotatebox[origin=c]{0}{Spécification} & \left\{
                    \begin{array}{p{.5cm}l}
                  \multicolumn{2}{l}{\mbox{Variables Pertinentes}} \\
            & \mbox{$
        %% 3 - variable pertinente
$} \\
                  \multicolumn{1}{l}{\mbox{Décomposition}}\\
            & \lefteqn{\mbox{$
        %% 4 - Conjointe
$}  }\\
            & \quad \mbox{$
        %% 5 - Decomposition
$} \\
            \multicolumn{2}{l}{\mbox{Formes Paramétriques}}\\
            & \mbox{ $
        %% 6 - Formes parametriques
$} \\
            \end{array} \right. \\
        \multicolumn{1}{l}{\mbox{Identification }}\\
                        & \mbox{$
        %% 7 - Identification
$} \\
        \end{array} \right. \\
    \multicolumn{1}{l}{\mbox{Question}} \\
                & \mbox{$ 
        %% 8 - Question
$}
    \end{array} \right.
\end{array}
\end{displaymath}}
\caption {Titre}
\label{label}
\end{figure}
\floatstyle{plain}
\restylefloat{figure}

}
%
\providecommand{\PBR}[8]{


\floatstyle{boxed}


\restylefloat{figure}
\begin{figure}[htbp]
{\small \begin{displaymath}
\begin{array}{ll}
    \rotatebox[origin=c]{90}{Programme} \left\{
          \begin{array}{ll}
                \rotatebox[origin=c]{0}{Description} & \left\{
        \begin{array}{ll}
        \rotatebox[origin=c]{0}{Spécification} & \left\{
                    \begin{array}{p{.5cm}l}
                  \multicolumn{2}{l}{\mbox{Variables Pertinentes}} \\
            & \mbox{#3} \\
                  \multicolumn{1}{l}{\mbox{Décomposition}}\\
            & \lefteqn{\mbox{#4}  }\\
            & \quad \mbox{#5} \\
            \multicolumn{2}{l}{\mbox{Formes Paramétriques}}\\
            & \mbox{ #6} \\
            \end{array} \right. \\
        \multicolumn{1}{l}{\mbox{Identification }}\\
                        & \mbox{#7} \\
        \end{array} \right. \\
    \multicolumn{1}{l}{\mbox{Question}} \\
                & \mbox{#8}
    \end{array} \right.
\end{array}
\end{displaymath}}
\caption {#1}
\label{#2}
\end{figure}
\floatstyle{plain}
\restylefloat{figure}
}
%
\providecommand{\PBRtest}[8]{


 
\floatstyle{boxed}

 
\restylefloat{figure}
\begin{figure}[htbp]
{\small \begin{displaymath}
\begin{array}{ll}
    \rotatebox[origin=c]{90}{Programme} \left\{
          \begin{array}{ll}
                \rotatebox[origin=c]{0}{Description} & \left\{
        \begin{array}{ll}
        \rotatebox[origin=c]{0}{Spécification} & \left\{
                    \begin{array}{p{.5cm}l}
                  \multicolumn{2}{l}{\mbox{Variables Pertinentes}} \\
            & \mbox{
        #3
} \\
                  \multicolumn{1}{l}{\mbox{Décomposition}}\\
            & \lefteqn{\mbox{
        #4
}  }\\
            & \quad \mbox{
        #5
} \\
            \multicolumn{2}{l}{\mbox{Formes Paramétriques}}\\
            & \mbox{ 
        #6
} \\
            \end{array} \right. \\
        \multicolumn{1}{l}{\mbox{Identification }}\\
                        & \mbox{
        #7
} \\
        \end{array} \right. \\
    \multicolumn{1}{l}{\mbox{Question}} \\
                & \mbox{
        #8
}
    \end{array} \right.
\end{array}
\end{displaymath}}
\caption { #1 }
\label{ #2 }
\end{figure}
\floatstyle{plain}
\restylefloat{figure}

}%
\providecommand{\PR}[2]{P(#1\:\mbox{\large \boldmath$|$}\:#2)}
\setpapersize{A4}
\setmarginsrb{1.5in}{1in}{1.00in}{1.00in}{14pt}{0.3in}{14pt}{0.494in}

\graphicspath{
  {./Figures/}
}
\let\footenote=\endnote

\fancypagestyle{myPageStyle}{
  \fancyhf{}
  \fancyhead[LO]{\normalsize\textit{\nouppercase{\rightmark}}}
  \fancyhead[RE]{\normalsize\textit{\nouppercase{\leftmark}}}
  \fancyhead[RO,LE]{\thepage }
  %
\renewcommand{\headrulewidth}{0.4pt}  %
\renewcommand{\footrulewidth}{0.4pt}}

\newlength{\larg}\setlength{\larg}{14.5cm}
\title{
	\begin{center}
		%%{\rule{\larg}{1mm}}
		
\vspace{4cm}
		Enfin en finir avec les problemes :\\
		{\bf Une mère a deux enfants dont l'un est un garçon. Quelle est la probabilité que l'autre soit une fille ?}\\
		et\\
		{\bf Le paradoxe des 3 portes}\\
		%%{\rule{\larg}{1mm}}
		
\vspace{8cm}
	\end{center}
}

\author{
	\begin{tabular}{p{13.7cm} l }  
		{\bf Note Technique }\\
		Pierre Dangauthier\\
		Projet EMotion\\
		\today
		\vspace{1mm}
	\end{tabular}\\
\hline
}\date{}

\setcounter{secnumdepth}{2}\setcounter{tocdepth}{1}
%
\providecommand{\etal}{{\it ~et al.~}}%
\providecommand{\eg}{{\em e.g.~}}%
\providecommand{\ie}{{\em i.e.~}}%
\providecommand{\etc}{{\em etc.~}}
%
\providecommand{\bs}{$\bigstar$}
%
\providecommand{\ed}{\;\;}%
\providecommand{\sd}{\;}
%
\providecommand{\pb}{programme bayésien }%
\providecommand{\pbs}{programmes bayésiens }%
\providecommand{\p}{probabilité }
%
\providecommand{\ebs}[1]{\mbox{\boldmath$#1$}}
%
\providecommand{\els}[3]{{^{#2}_{#3}#1}}%
\providecommand{\ers}[3]{{#1^{#2}_{#3}}}%
\providecommand{\elrs}[5]{{^{#2}_{#3}#1^{#4}_{#5}}}
%
\providecommand{\ecl}[1]{\mathcal{#1}}


\usepackage[dvips]{color}


\pagecolor[gray]{.7}

\usepackage[latin1]{inputenc}



\makeatletter

\makeatletter
\count@=\the\catcode`\_ \catcode`\_=8 
\newenvironment{tex2html_wrap}{}{}%
\catcode`\<=12\catcode`\_=\count@
\newcommand{\providedcommand}[1]{\expandafter\providecommand\csname #1\endcsname}%
\newcommand{\renewedcommand}[1]{\expandafter\providecommand\csname #1\endcsname{}%
  \expandafter\renewcommand\csname #1\endcsname}%
\newcommand{\newedenvironment}[1]{\newenvironment{#1}{}{}\renewenvironment{#1}}%
\let\newedcommand\renewedcommand
\let\renewedenvironment\newedenvironment
\makeatother
\let\mathon=$
\let\mathoff=$
\ifx\AtBeginDocument\undefined \newcommand{\AtBeginDocument}[1]{}\fi
\newbox\sizebox
\setlength{\hoffset}{0pt}\setlength{\voffset}{0pt}
\addtolength{\textheight}{\footskip}\setlength{\footskip}{0pt}
\addtolength{\textheight}{\topmargin}\setlength{\topmargin}{0pt}
\addtolength{\textheight}{\headheight}\setlength{\headheight}{0pt}
\addtolength{\textheight}{\headsep}\setlength{\headsep}{0pt}
\setlength{\textwidth}{349pt}
\newwrite\lthtmlwrite
\makeatletter
\let\realnormalsize=\normalsize
\global\topskip=2sp
\def\preveqno{}\let\real@float=\@float \let\realend@float=\end@float
\def\@float{\let\@savefreelist\@freelist\real@float}
\def\liih@math{\ifmmode$\else\bad@math\fi}
\def\end@float{\realend@float\global\let\@freelist\@savefreelist}
\let\real@dbflt=\@dbflt \let\end@dblfloat=\end@float
\let\@largefloatcheck=\relax
\let\if@boxedmulticols=\iftrue
\def\@dbflt{\let\@savefreelist\@freelist\real@dbflt}
\def\adjustnormalsize{\def\normalsize{\mathsurround=0pt \realnormalsize
 \parindent=0pt\abovedisplayskip=0pt\belowdisplayskip=0pt}%
 \def\phantompar{\csname par\endcsname}\normalsize}%
\def\lthtmltypeout#1{{\let\protect\string \immediate\write\lthtmlwrite{#1}}}%
\newcommand\lthtmlhboxmathA{\adjustnormalsize\setbox\sizebox=\hbox\bgroup\kern.05em }%
\newcommand\lthtmlhboxmathB{\adjustnormalsize\setbox\sizebox=\hbox to\hsize\bgroup\hfill }%
\newcommand\lthtmlvboxmathA{\adjustnormalsize\setbox\sizebox=\vbox\bgroup %
 \let\ifinner=\iffalse \let\)\liih@math }%
\newcommand\lthtmlboxmathZ{\@next\next\@currlist{}{\def\next{\voidb@x}}%
 \expandafter\box\next\egroup}%
\newcommand\lthtmlmathtype[1]{\gdef\lthtmlmathenv{#1}}%
\newcommand\lthtmllogmath{\lthtmltypeout{l2hSize %
:\lthtmlmathenv:\the\ht\sizebox::\the\dp\sizebox::\the\wd\sizebox.\preveqno}}%
\newcommand\lthtmlfigureA[1]{\let\@savefreelist\@freelist
       \lthtmlmathtype{#1}\lthtmlvboxmathA}%
\newcommand\lthtmlpictureA{\bgroup\catcode`\_=8 \lthtmlpictureB}%
\newcommand\lthtmlpictureB[1]{\lthtmlmathtype{#1}\egroup
       \let\@savefreelist\@freelist \lthtmlhboxmathB}%
\newcommand\lthtmlpictureZ[1]{\hfill\lthtmlfigureZ}%
\newcommand\lthtmlfigureZ{\lthtmlboxmathZ\lthtmllogmath\copy\sizebox
       \global\let\@freelist\@savefreelist}%
\newcommand\lthtmldisplayA{\bgroup\catcode`\_=8 \lthtmldisplayAi}%
\newcommand\lthtmldisplayAi[1]{\lthtmlmathtype{#1}\egroup\lthtmlvboxmathA}%
\newcommand\lthtmldisplayB[1]{\edef\preveqno{(\theequation)}%
  \lthtmldisplayA{#1}\let\@eqnnum\relax}%
\newcommand\lthtmldisplayZ{\lthtmlboxmathZ\lthtmllogmath\lthtmlsetmath}%
\newcommand\lthtmlinlinemathA{\bgroup\catcode`\_=8 \lthtmlinlinemathB}
\newcommand\lthtmlinlinemathB[1]{\lthtmlmathtype{#1}\egroup\lthtmlhboxmathA
  \vrule height1.5ex width0pt }%
\newcommand\lthtmlinlineA{\bgroup\catcode`\_=8 \lthtmlinlineB}%
\newcommand\lthtmlinlineB[1]{\lthtmlmathtype{#1}\egroup\lthtmlhboxmathA}%
\newcommand\lthtmlinlineZ{\egroup\expandafter\ifdim\dp\sizebox>0pt %
  \expandafter\centerinlinemath\fi\lthtmllogmath\lthtmlsetinline}
\newcommand\lthtmlinlinemathZ{\egroup\expandafter\ifdim\dp\sizebox>0pt %
  \expandafter\centerinlinemath\fi\lthtmllogmath\lthtmlsetmath}
\newcommand\lthtmlindisplaymathZ{\egroup %
  \centerinlinemath\lthtmllogmath\lthtmlsetmath}
\def\lthtmlsetinline{\hbox{\vrule width.1em \vtop{\vbox{%
  \kern.1em\copy\sizebox}\ifdim\dp\sizebox>0pt\kern.1em\else\kern.3pt\fi
  \ifdim\hsize>\wd\sizebox \hrule depth1pt\fi}}}
\def\lthtmlsetmath{\hbox{\vrule width.1em\kern-.05em\vtop{\vbox{%
  \kern.1em\kern0.8 pt\hbox{\hglue.17em\copy\sizebox\hglue0.8 pt}}\kern.3pt%
  \ifdim\dp\sizebox>0pt\kern.1em\fi \kern0.8 pt%
  \ifdim\hsize>\wd\sizebox \hrule depth1pt\fi}}}
\def\centerinlinemath{%
  \dimen1=\ifdim\ht\sizebox<\dp\sizebox \dp\sizebox\else\ht\sizebox\fi
  \advance\dimen1by.5pt \vrule width0pt height\dimen1 depth\dimen1 
 \dp\sizebox=\dimen1\ht\sizebox=\dimen1\relax}

\def\lthtmlcheckvsize{\ifdim\ht\sizebox<\vsize 
  \ifdim\wd\sizebox<\hsize\expandafter\hfill\fi \expandafter\vfill
  \else\expandafter\vss\fi}%
\providecommand{\selectlanguage}[1]{}%
\makeatletter \tracingstats = 1 
\providecommand{\Beta}{\textrm{B}}
\providecommand{\Mu}{\textrm{M}}
\providecommand{\Kappa}{\textrm{K}}
\providecommand{\Rho}{\textrm{R}}
\providecommand{\Epsilon}{\textrm{E}}
\providecommand{\Chi}{\textrm{X}}
\providecommand{\Iota}{\textrm{J}}
\providecommand{\omicron}{\textrm{o}}
\providecommand{\Zeta}{\textrm{Z}}
\providecommand{\Eta}{\textrm{H}}
\providecommand{\Omicron}{\textrm{O}}
\providecommand{\Nu}{\textrm{N}}
\providecommand{\Tau}{\textrm{T}}
\providecommand{\Alpha}{\textrm{A}}


\begin{document}
\pagestyle{empty}\thispagestyle{empty}\lthtmltypeout{}%
\lthtmltypeout{latex2htmlLength hsize=\the\hsize}\lthtmltypeout{}%
\lthtmltypeout{latex2htmlLength vsize=\the\vsize}\lthtmltypeout{}%
\lthtmltypeout{latex2htmlLength hoffset=\the\hoffset}\lthtmltypeout{}%
\lthtmltypeout{latex2htmlLength voffset=\the\voffset}\lthtmltypeout{}%
\lthtmltypeout{latex2htmlLength topmargin=\the\topmargin}\lthtmltypeout{}%
\lthtmltypeout{latex2htmlLength topskip=\the\topskip}\lthtmltypeout{}%
\lthtmltypeout{latex2htmlLength headheight=\the\headheight}\lthtmltypeout{}%
\lthtmltypeout{latex2htmlLength headsep=\the\headsep}\lthtmltypeout{}%
\lthtmltypeout{latex2htmlLength parskip=\the\parskip}\lthtmltypeout{}%
\lthtmltypeout{latex2htmlLength oddsidemargin=\the\oddsidemargin}\lthtmltypeout{}%
\makeatletter
\if@twoside\lthtmltypeout{latex2htmlLength evensidemargin=\the\evensidemargin}%
\else\lthtmltypeout{latex2htmlLength evensidemargin=\the\oddsidemargin}\fi%
\lthtmltypeout{}%
\makeatother
\setcounter{page}{1}
\onecolumn

% !!! IMAGES START HERE !!!

\setcounter{secnumdepth}{2}
\setcounter{tocdepth}{1}
\stepcounter{section}
{\newpage\clearpage
\lthtmlinlinemathA{tex2html_wrap_inline1204}%
$ \frac{1}{2}$%
\lthtmlinlinemathZ
\lthtmlcheckvsize\clearpage}

{\newpage\clearpage
\lthtmlinlinemathA{tex2html_wrap_inline1206}%
$ \frac{2}{3}$%
\lthtmlinlinemathZ
\lthtmlcheckvsize\clearpage}

\stepcounter{section}
{\newpage\clearpage
\lthtmlinlinemathA{tex2html_wrap_inline1240}%
$ N_1$%
\lthtmlinlinemathZ
\lthtmlcheckvsize\clearpage}

{\newpage\clearpage
\lthtmlinlinemathA{tex2html_wrap_inline1242}%
$ N_2$%
\lthtmlinlinemathZ
\lthtmlcheckvsize\clearpage}

{\newpage\clearpage
\lthtmlinlinemathA{tex2html_wrap_inline1244}%
$ n_f$%
\lthtmlinlinemathZ
\lthtmlcheckvsize\clearpage}

{\newpage\clearpage
\lthtmlinlinemathA{tex2html_wrap_inline1246}%
$ \frac{n_f}{N_2}$%
\lthtmlinlinemathZ
\lthtmlcheckvsize\clearpage}

\stepcounter{section}
\stepcounter{section}
\stepcounter{subsection}
\stepcounter{subsection}
\stepcounter{subsection}
\stepcounter{section}
\stepcounter{subsection}
{\newpage\clearpage
\lthtmlinlinemathA{tex2html_wrap_inline1329}%
$ P(ff | I=1)=0$%
\lthtmlinlinemathZ
\lthtmlcheckvsize\clearpage}

{\newpage\clearpage
\lthtmlinlinemathA{tex2html_wrap_inline1331}%
$ P(gf | I=1)=...=1/3$%
\lthtmlinlinemathZ
\lthtmlcheckvsize\clearpage}

\stepcounter{subsection}
{\newpage\clearpage
\lthtmlinlinemathA{tex2html_wrap_inline1377}%
$ P(A|B)$%
\lthtmlinlinemathZ
\lthtmlcheckvsize\clearpage}

{\newpage\clearpage
\lthtmlinlinemathA{tex2html_wrap_indisplay1379}%
$\displaystyle P(A|B) = P($%
\lthtmlindisplaymathZ
\lthtmlcheckvsize\clearpage}

{\newpage\clearpage
\lthtmlinlinemathA{tex2html_wrap_indisplay1380}%
$\displaystyle ) = \frac{P(\text{A et B})}{P(B)}$%
\lthtmlindisplaymathZ
\lthtmlcheckvsize\clearpage}

{\newpage\clearpage
\lthtmlinlinemathA{tex2html_wrap_inline1382}%
$ \frac{1}{4}$%
\lthtmlinlinemathZ
\lthtmlcheckvsize\clearpage}

{\newpage\clearpage
\lthtmlinlinemathA{tex2html_wrap_inline1386}%
$ P($%
\lthtmlinlinemathZ
\lthtmlcheckvsize\clearpage}

{\newpage\clearpage
\lthtmlinlinemathA{tex2html_wrap_inline1387}%
$ ) =\frac{ P(\text{ 2G et "il y a un g"})}{P(\text{"il y a un g"})} = \frac{P( \text{2G} )}{ P(\text{"il y a un g"})} = \frac{1/4}{3/4} = \frac{1}{3}$%
\lthtmlinlinemathZ
\lthtmlcheckvsize\clearpage}

{\newpage\clearpage
\lthtmlinlinemathA{tex2html_wrap_inline1390}%
$ ) =\frac{ P(\text{ 2F et "il y a un g"}) }{ P(\text{"il y a un g"}}  = \frac{P ( \text{impossible}) }{ P(\text{"il y a un g"}}  = \frac{0}{ 3/4} = 0$%
\lthtmlinlinemathZ
\lthtmlcheckvsize\clearpage}

{\newpage\clearpage
\lthtmlinlinemathA{tex2html_wrap_inline1393}%
$ ) =\frac{ P( \text{Mix et "il y a un g"}) }{ P(\text{"il y a un g"})}  = \frac{P ( \text{Mix}) }{ P(\text{"il y a un g"})}  = \frac{1/2}{3/4} =  \frac{2}{3}$%
\lthtmlinlinemathZ
\lthtmlcheckvsize\clearpage}

{\newpage\clearpage
\lthtmlinlinemathA{tex2html_wrap_inline1395}%
$ F \in \{0,1,2\}$%
\lthtmlinlinemathZ
\lthtmlcheckvsize\clearpage}

{\newpage\clearpage
\lthtmlinlinemathA{tex2html_wrap_inline1397}%
$ F$%
\lthtmlinlinemathZ
\lthtmlcheckvsize\clearpage}

{\newpage\clearpage
\lthtmlfigureA{figure553}%
\begin{figure}{\small\begin{displaymath}
\begin{array}{ll}
    \rotatebox[origin=c]{90}{Programme}\left\{
          \begin{array}{ll}
                \rotatebox[origin=c]{0}{Description}& \left\{
        \begin{array}{ll}
        \rotatebox[origin=c]{0}{Spécification}& \left\{
                    \begin{array}{p{.5cm}l}
                  \multicolumn{2}{l}{\mbox{Variables Pertinentes}} \\
            & \mbox{F} \\
                  \multicolumn{1}{l}{\mbox{Décomposition}}\\
            & \lefteqn{\mbox{$P(F)$}  }\\
            & \quad \mbox{} \\
            \multicolumn{2}{l}{\mbox{Formes Paramétriques}}\\
            & \mbox{ Histogramme} \\
            \end{array} \right. \\
        \multicolumn{1}{l}{\mbox{Identification }}\\
                        & \mbox{$P(F=0)=P(F=2)=1/4$ et $P(F=1)=1/2$} \\
        \end{array} \right. \\
    \multicolumn{1}{l}{\mbox{Question}} \\
                & \mbox{$P( F=1 | F \not = 2) = \frac{P(F \not = 2 | F=1 ) P(F=1)}{p(F \not = 2)} = \frac{  1 * 0.5}{1-0.25}= 2/3$}
    \end{array} \right.
\end{array}
\end{displaymath}}
\end{figure}%
\lthtmlfigureZ
\lthtmlcheckvsize\clearpage}

\stepcounter{section}
\stepcounter{section}
\stepcounter{subsection}
\stepcounter{subsection}
\stepcounter{subsection}
\stepcounter{subsection}
{\newpage\clearpage
\lthtmlinlinemathA{tex2html_wrap_inline1498}%
$ \frac{1}{3}$%
\lthtmlinlinemathZ
\lthtmlcheckvsize\clearpage}

\stepcounter{subsection}
{\newpage\clearpage
\lthtmlinlinemathA{tex2html_wrap_inline1528}%
$ H_i$%
\lthtmlinlinemathZ
\lthtmlcheckvsize\clearpage}

{\newpage\clearpage
\lthtmlinlinemathA{tex2html_wrap_indisplay1532}%
$\displaystyle P(H_1)=P(H_2)=P(H_3)=\frac{1}{3}.$%
\lthtmlindisplaymathZ
\lthtmlcheckvsize\clearpage}

{\newpage\clearpage
\lthtmlinlinemathA{tex2html_wrap_inline1534}%
$ O$%
\lthtmlinlinemathZ
\lthtmlcheckvsize\clearpage}

{\newpage\clearpage
\lthtmlinlinemathA{tex2html_wrap_inline1536}%
$ P( O=2 | H_1) = \frac{1}{2}$%
\lthtmlinlinemathZ
\lthtmlcheckvsize\clearpage}

{\newpage\clearpage
\lthtmlinlinemathA{tex2html_wrap_inline1538}%
$ P( O=3 | H_1) = \frac{1}{2}$%
\lthtmlinlinemathZ
\lthtmlcheckvsize\clearpage}

{\newpage\clearpage
\lthtmlinlinemathA{tex2html_wrap_inline1540}%
$ P( O=2 | H_2) = 0$%
\lthtmlinlinemathZ
\lthtmlcheckvsize\clearpage}

{\newpage\clearpage
\lthtmlinlinemathA{tex2html_wrap_inline1542}%
$ P( O=3 | H_2) = 1$%
\lthtmlinlinemathZ
\lthtmlcheckvsize\clearpage}

{\newpage\clearpage
\lthtmlinlinemathA{tex2html_wrap_inline1544}%
$ P( O=2 | H_3) = 1$%
\lthtmlinlinemathZ
\lthtmlcheckvsize\clearpage}

{\newpage\clearpage
\lthtmlinlinemathA{tex2html_wrap_inline1546}%
$ P( O=3 | H_3) = 0$%
\lthtmlinlinemathZ
\lthtmlcheckvsize\clearpage}

{\newpage\clearpage
\lthtmlinlinemathA{tex2html_wrap_inline1548}%
$ O=3$%
\lthtmlinlinemathZ
\lthtmlcheckvsize\clearpage}

{\newpage\clearpage
\lthtmlinlinemathA{tex2html_wrap_inline1550}%
$ i$%
\lthtmlinlinemathZ
\lthtmlcheckvsize\clearpage}

{\newpage\clearpage
\lthtmlinlinemathA{tex2html_wrap_indisplay1552}%
$\displaystyle P(H_i | O=3) = \frac{P(O=3 | H_i) P(H_i)}{P(O=3)}$%
\lthtmlindisplaymathZ
\lthtmlcheckvsize\clearpage}

{\newpage\clearpage
\lthtmlinlinemathA{tex2html_wrap_inline1554}%
$ P(H_i)$%
\lthtmlinlinemathZ
\lthtmlcheckvsize\clearpage}

{\newpage\clearpage
\lthtmlinlinemathA{tex2html_wrap_inline1556}%
$ P(O=3)$%
\lthtmlinlinemathZ
\lthtmlcheckvsize\clearpage}

{\newpage\clearpage
\lthtmlinlinemathA{tex2html_wrap_inline1562}%
$ P(O=3 | H_i)$%
\lthtmlinlinemathZ
\lthtmlcheckvsize\clearpage}

{\newpage\clearpage
\lthtmlinlinemathA{tex2html_wrap_inline1564}%
$ \frac{1}{2},1,0$%
\lthtmlinlinemathZ
\lthtmlcheckvsize\clearpage}

{\newpage\clearpage
\lthtmlinlinemathA{tex2html_wrap_inline1566}%
$ i=1,2,3$%
\lthtmlinlinemathZ
\lthtmlcheckvsize\clearpage}

{\newpage\clearpage
\lthtmlinlinemathA{tex2html_wrap_inline1568}%
$ H_2$%
\lthtmlinlinemathZ
\lthtmlcheckvsize\clearpage}

{\newpage\clearpage
\lthtmlinlinemathA{tex2html_wrap_inline1570}%
$ H_1$%
\lthtmlinlinemathZ
\lthtmlcheckvsize\clearpage}

{\newpage\clearpage
\lthtmlinlinemathA{tex2html_wrap_indisplay1574}%
$\displaystyle P(O=3)=\sum_i{P( O=3 | H_i) P(H_i)}=\frac{1/2 + 1 + 0}{3}=\frac{1}{2}$%
\lthtmlindisplaymathZ
\lthtmlcheckvsize\clearpage}

{\newpage\clearpage
\lthtmlinlinemathA{tex2html_wrap_indisplay1576}%
$\displaystyle P(H_1 | O=3) = \frac{1}{3} \: ; \:  P(H_2 | O=3) = \frac{2}{3} \: ; \: P(H_3 | O=3) = 0 $%
\lthtmlindisplaymathZ
\lthtmlcheckvsize\clearpage}


\end{document}
